% Drawing 29/10-3 -- folder 29, batch 10, archive page 183 (PDF page 184).
% Transcription: moitié gauche de la feuille, six petites esquisses à l'encre : un triangle très effilé fait d'une flèche épaisse descendant vers la gauche et d'une flèche fine revenant vers la droite, avec une flèche courbe entre leurs deux bouts ; un petit arc issu d'un point ; un point $x$ relié par une flèche à un point $a$ qui porte une boucle ; une chaîne de deux flèches, la première marquée $u_1$, d'où partent au bout deux flèches divergentes ; trois ovales d'où descendent des flèches ; une ligne brisée fléchée issue d'un point.
% Kind: six small sketches of arrows, loops and paths. Method 2 (schematic redraw from the facsimile). Status: partial.
% Notes: (a) The thin triangle: a heavy arrow from a marked point down to the lower left, a thin arrow from its tip back to the right, a second marked point just beyond that arrow's head and a heavy curved arrow from it up to the first point. (b) A small arc starting at a small ring. (c) The point x, an arrow to the point a, both lettered, and a loop (a petal) on a with its arrowhead coming back to a. (d) A marked point, an arrow to a second marked point (u_1 written above it), a heavy arrow from there to a node, and from the node two arrows, one up to the right, one to the right; a small arrow beside them. (e) Three ovals: two arrows go down from the first, a stroke inside the second, a heavy arrow down to the left from the third. (f) A marked point and a broken line of three segments ending in an arrowhead. Overlaid on the 400 dpi crop. Still differs: the three ovals are free loops drawn in one stroke each, here tilted ellipses; the small arrow beside the fan of (d) is hooked on the page. Not drawn: the letters E, K_1, f(x,y), f(u) (text), the vertical stroke and dot under E, the long thin oblique line that crosses the sheet. Redrawn 2026-11-12 (issue 128): to check again.
% Coordinates: 1 unit = 1 mm = 1/100 of the crop width (bbox in page fractions: [0.1, 0.27, 0.48, 0.53]).
% Redrawn by Claude (claude-fable-5-2), issue #128, 2026-11-12; not checked by a human; not compiled (no TeX in the session).
% Rights: reproduces part of the work -- behind RENDER_TRANSCRIPTS, never published by itself.
\documentclass[tikz,border=3pt]{standalone}
\usepackage{amsmath,amssymb}
\usetikzlibrary{patterns,arrows.meta}
\tikzset{pencil/.style={black!45}, struck/.style={black!35}, mpt/.style={circle,draw,inner sep=1.3pt}, pt/.style={circle,fill,inner sep=0.9pt}, opt/.style={circle,draw,fill=white,inner sep=1.1pt}}
\begin{document}
\begin{tikzpicture}[x=1mm,y=1mm,line cap=round,line join=round,font=\small]
\node[pt] at (24.6,-19.6) {};
\draw[-{Stealth[length=1.8mm]}, very thick] (24.6,-19.6) -- (7,-45);
\draw[-{Stealth[length=1.8mm]}] (7,-45) -- (25.6,-34.2);
\node[pt] at (27.6,-32.6) {};
\draw[-{Stealth[length=1.8mm]}, very thick] plot[smooth, tension=0.55] coordinates {(27.6,-32.6) (30,-29) (30.2,-24) (27,-19.8)};
\node[opt] at (70.4,-11.6) {};
\draw plot[smooth, tension=0.55] coordinates {(70.8,-11.6) (75.6,-10) (79.2,-5.4)};
\node[pt] at (59.6,-33) {};
\node[pt] at (72.4,-27.4) {};
\draw[-{Stealth[length=1.8mm]}] (59.6,-33) -- (72.4,-27.4);
\node at (60,-36) {$x$};
\node at (73,-30.4) {$a$};
\draw[-{Stealth[length=1.8mm]}] plot[smooth, tension=0.8] coordinates {(73.8,-25.6) (77.2,-19.5) (76.6,-15.6) (74.6,-16.5) (73,-21) (72.4,-25.8)};
\node[pt] at (47.6,-60.4) {};
\draw[-{Stealth[length=1.8mm]}, thick] (47.6,-60.4) -- (62.4,-56.8);
\node[pt] at (63.4,-56.6) {};
\draw[-{Stealth[length=1.8mm]}, very thick] (63.4,-56.6) -- (74.6,-49.6);
\draw[-{Stealth[length=1.8mm]}, thick] (75.2,-48.8) -- (80.8,-38.2);
\draw[-{Stealth[length=1.8mm]}, very thick] (75.2,-50.2) -- (86.4,-50.4);
\draw[-{Stealth[length=1.8mm]}] (84.4,-46) -- (82.8,-42.6);
\node at (56,-51.6) {$u_1$};
\draw[rotate around={-15:(76,-68.3)}] (76,-68.3) ellipse[x radius=3, y radius=6.2];
\draw[-{Stealth[length=1.8mm]}] (74.2,-71) -- (75.8,-84.6);
\draw[-{Stealth[length=1.8mm]}] (79.2,-72) -- (77.6,-84.4);
\draw[rotate around={-13:(83,-70.3)}] (83,-70.3) ellipse[x radius=2.8, y radius=7];
\draw (83.6,-71) -- (81.4,-77);
\draw[rotate around={-22:(89.6,-76)}] (89.6,-76) ellipse[x radius=3, y radius=6.8];
\draw[-{Stealth[length=1.8mm]}, thick] (90.4,-74) -- (80.6,-86);
\node[pt] at (41,-91) {};
\draw[-{Stealth[length=1.8mm]}] (41,-91) -- (48,-86.4) -- (54.4,-90.8) -- (61,-82.6);
\end{tikzpicture}
\end{document}
